\section{Лекция 22}
% \begin{statement}
%     Площадь поверхности сферы вписанной в цилиндр, равна площади боковой поверхности этого цилиндра.
% \end{statement}
% \begin{proof}
%     \text{Нет в экзамене, но, возможно, появится позже.}
% \end{proof}
% Тут картинка цилиндр описаный вокруг сферы
% Берем точку $x$ на сфере и отображаем ее в точку на поверхности цилиндра, проведя $x$ перпендикуляр к оси цилиндра. Хотим доказать, что площадь поверхности сферы будет равна площади поверхности цилиндра. По теореме Гульдина площадь кусочка сферы между $a$ и $b$ равна
% \[S_{a,b}=|\gamma|\cdot 2\pi l\]
% при этом площадь 
% тут либо некст теорема следует из этого
% лиюо это следует из некст теоремы предельным переходом
\begin{theorem} [Теорема Пуассона] $\\$
    Пусть $f(x,y,z)=f(z)$ (зависит только от $z$) определена и непрерывна на сфере 
    \[x^2+y^2+z^2=R^2\]
    Тогда
    \[\iint\limits_{x^2+y^2+z^2=R^2}f(z)\ dS=2\pi R\cdot \int\limits_{-R}^{R}f(z)\ dz\]
\end{theorem}
\begin{proof}
    Параметризуем сферу с использованием цилиндрических координат
    \[\begin{cases}
        x=r\cdot\cos{\phi},\\
        y=r\cdot \sin{\phi},\\
        z=z
    \end{cases}\]
    При этом мы находимся на сфере
    \[x^2+y^2+z^2=R^2\ \Rightarrow\ r^2+z^2=R^2\ \Rightarrow\ r=\sqrt{R^2-z^2}\]
    Итак, искомая параметризация    
    \[\begin{cases}
        x=\sqrt{R^2-z^2}\cdot\cos{\phi},\\
        y=\sqrt{R^2-z^2}\cdot \sin{\phi},\\
        z=z 
    \end{cases}\]
    где $z\in[-R,R],\ \phi\in [0,2\pi]$. Значит
    \[r'_z=\left(\frac{-z}{\sqrt{R^2-z^2}}\cdot \cos{\phi},\ \frac{-z}{\sqrt{R^2-z^2}}\sin{\phi},\ 1\right)\]
    \[r'_{\phi}=\left(-\sqrt{R^2-z^2}\cdot \sin{\phi},\ \sqrt{R^2-z^2}\cdot \cos{\phi},\ 0\right)\]
    Таким образом
    \[E=\frac{z^2}{R^2-z^2}\cdot (\cos^2{\phi}+\sin^2{\phi})+1=\frac{R^2}{R^2-z^2}\]
    \[G=(R^2-z^2)(\sin^2{\phi}+\cos^2{\phi})=R^2-z^2\] 
    \[F=\frac{z}{\sqrt{R^2-z^2}}\cdot \cos{\phi}\cdot \sqrt{R^2-z^2}\cdot\sin{\phi}-\frac{z}{\sqrt{R^2-z^2}}\cdot\sin{\phi}\cdot \sqrt{R^2-z^2}\cdot\cos{\phi}=0\]
    Отсюда
    \[\sqrt{EG-F^2}=\sqrt{\frac{R^2}{R^2-z^2}\cdot (R^2-z^2)}=R\]
    тогда 
    \[\iint\limits_{x^2+y^2+z^2=R^2}f(z)\ dS=\int\limits_{0}^{2\pi}\int\limits_{-R}^{R}f(z)\cdot R\ dzd\phi=2\pi R\cdot \int\limits_{-R}^{R}f(z)\ dz\]
\end{proof}
\noindentВ этой лекции еще было утверждение про то, что площадь поверхности сферы, вписанной в цилиндр, равна площади боковой поверхности этого цилиндра. Также была обсуждена стереографическая проекция. (\textit{этого нет в программе экзамена})
% тут что то про стереографическую и азимутальную проекции (Косухин сказал картинки можно посмотреть на сайте математический этюд)
\newpage